\section{Crossed products and duality\label{sec:NAdef}}

In this section we summarise some of the mathematical tools that we
will use in this paper. A~good reference for the material covered in
the following is the book~\cite{Williams2007}. Throughout this paper,
all topological spaces are assumed to be second countable (hence
separable), locally compact and Hausdorff.

\subsection{Dynamical systems and their crossed products}

Let $\sfG$ be a locally compact group and let $X$ be a $\sfG$-space, i.e., a
topological space which is acted
upon
homeomorphically by $\sfG$; we denote the $\sfG$-action $\sfG\times
X\to X$ by $(\gamma,x)\mapsto\gamma\cdot x$. The pair $(X,\sfG)$ is
called a \emph{transformation group}.
 A related concept is that of a
\emph{dynamical system}, which is a triple $(\alg,\sfG,\alpha)$
consisting of an algebra $\alg$ and a locally compact group $\sfG$ acting
on~$\alg$ via a group homomorphism $\alpha\colon \sfG\to\Aut (\alg)$,
denoted $\gamma\mapsto(\alpha_\gamma\colon \alg\to\alg)$ for $\gamma\in
\sfG$. In topological
T-duality one usually requires $\alg$ to be a $C^*$-algebra, in which
case $(\alg,\sfG,\alpha)$ is called a \emph{$C^*$-dynamical
 system}. Two dynamical systems $(\alg,\sfG,\alpha)$ and $(\balg,\sfG,\beta)$ are
\emph{equivalent} if there is an algebra isomorphism
$\varphi\colon \alg\to\balg$ which intertwines the $\sfG$-actions: $\varphi\circ\alpha_\gamma =
\beta_\gamma\circ\varphi$ for all $\gamma\in \sfG$.

If $\alg$ is a commutative $C^*$-algebra, then we call $(\alg,\sfG,\alpha)$ a
\emph{commutative} dynamical system. In that case, by Gelfand duality
$\alg=C_0(X)$ is the algebra of continuous functions vanishing at
infinity on a
topological space $X$
equipped with the $\sfG$-action \smash{$\alpha_\gamma^\dag\big|_X$} for $\gamma\in \sfG$ under
the identification of points $x\in X$ with irreducible representations
of $C_0(X)$, which are one-dimensional and given by the point evaluation
maps ${\rm ev}_x\colon \alg\to\alg$ with ${\rm ev}_x(\mathtt{f})=\mathtt{f}(x)$; then
$(X,\sfG)$ is a~transformation group. Conversely, given a transformation
group $(X,\sfG)$, there is an associated commutative dynamical system
$(C_0(X), \sfG,\alpha)$, where $\alpha_\gamma(\mathtt{f})(x)=\mathtt{f}\big(\gamma^{-1}\cdot x\big)$ for $\gamma\in
\sfG$, $\mathtt{f}\in C_0(X)$ and $x\in X$. In other words, there is a one-to-one
correspondence between transformation groups and commutative $C^*$-dynamical
systems.

If the $C^*$-algebra $\alg$ is not commutative, then we call
$(\alg,\sfG,\alpha)$ a \emph{noncommutative} $C^*$-dynamical system.

As usual, it is more useful to work with a representation rather than the
abstract dynamical system itself. A \emph{covariant representation} of
a dynamical system $(\alg,\sfG,\alpha)$ in a $C^*$-algebra $\balg$ with
multiplier algebra ${\rm M}(\balg)$ is a
pair $(\Pim,U)$ consisting of a~homomorphism $\Pim\colon \alg\to {\rm M}(\balg)$ and a~unitary representation $U\colon \sfG\to{\rm M}(\balg)$, $\gamma\mapsto U_\gamma$ which satisfy the compatibility condition
\[
\Pim\big(\alpha_\gamma(a)\big) = U_\gamma \Pim(a) U_\gamma^{-1} ,
\]
for all $\gamma\in \sfG$ and $a\in\alg$. A natural choice is to take
$\balg=\CK(\hil)$ to be the $C^*$-algebra of compact operators on a
separable Hilbert space $\hil$, which gives a
representation $\Pim\colon \alg\to{\rm B}(\hil)$ of the algebra $\alg$ by
bounded operators ${\rm B}(\hil)$ on $\hil$ and a unitary representation
$U\colon \sfG\to{\rm B}(\hil)$ of the group $\sfG$ on~$\hil$; in this case
we call $(\Pim,U)$ a covariant representation of $(\alg,\sfG,\alpha)$ on~$\hil$.

When a group $\sfG$ acts on a space $X$, one is naturally interested
in considering the
quotient space $X/\sfG$ of $\sfG$-orbits on~$X$. When $\sfG$ acts
freely and properly on
$X$, this is described algebraically by the algebra of functions
$C_0(X/\sfG)$. More generally, the subalgebra of $\sfG$-invariant elements
$\alg^\sfG\subseteq\alg$ of a $\sfG$-algebra $\alg$ can be used to represent the quotient, even for
$\sfG$-actions with fixed points. A more general and systematic way of
dealing with the effective algebraic ``quotient'' is through the
\emph{crossed product} algebra $\alg\rtimes_\alpha \sfG$ for a dynamical
system $(\alg,\sfG,\alpha)$. For a transformation group $(X,\sfG)$,
this description is particularly powerful in the cases where the
quotient $X/\sfG$ is not a Hausdorff space, while for a free and proper
$\sfG$-action it gives an algebra with the same spectrum $X/\sfG$
as the algebra $C_0(X/\sfG)=C_0(X)^\sfG$ of
$\sfG$-invariant functions on $X$.

In order to define the crossed product
algebra, we first define
\[
\|f\|_{\rm univ} := \sup_{({\mit\Pi},U)} \big\|({\mit\Pi}\rtimes_\alpha U)(f)\big\|
\]
for compactly supported functions $f\in C_{\rm c}(\sfG,\alg) $, where the supremum is taken
over (possibly degenerate) covariant representations
$({\mit\Pi},U)$ of $(\alg, \sfG, \alpha)$ with
\[
(\Pim\rtimes_\alpha U)(f) := \int_\sfG
\Pim\big(f(\gamma)\big) U_\gamma\, \dd\mu_\sfG(\gamma) ,
\]
and $\mu_\sfG$ denotes the left invariant Haar measure on $\sfG$. This defines a norm, called the universal norm, on the
space $C_{\rm c}(\sfG,\alg)$.
Then the crossed product algebra $\alg\rtimes_\alpha \sfG$ is the
completion (in the universal norm) of the algebra $C_{\rm c}(\sfG,\alg)$
equipped with the convolution product
\begin{gather}\label{eq:convproduct}
(f\star f')(\gamma) := \int_\sfG\, f(\gamma') \alpha_{\gamma'}\big(f'\big(\gamma'^{-1} \gamma\big)\big)\, \dd \mu_\sfG(\gamma')
,
\end{gather}
for all $f,f'\colon \sfG\to\alg$. In general this is a noncommutative
multiplication, even for commutative dynamical systems.
Since $\alg$ is a $C^*$-algebra, there is a $*$-structure on the
convolution algebra defined by
\[
f^\dag(\gamma):= \Delta_\sfG(\gamma)^{-1} \alpha_\gamma\big(f\big(\gamma^{-1}\big)^*\big) ,
\]
where $\Delta_\sfG\colon \sfG\to \real^+$ is the modular function of
$\sfG$ defined through
\[
\Delta_\sfG(\gamma')
\int_\sfG\,f(\gamma)\,\dd\mu_\sfG(\gamma)=\int_\sfG f(\gamma \gamma')\,\dd\mu_\sfG(\gamma)
\]
for $f\in C_{\rm c}(\sfG,\alg)$ and $\gamma'\in \sfG$. By the uniqueness of the left
invariant Haar measure $\mu_\sfG$ up to a~positive
constant, $\Delta_\sfG(\gamma')$ is independent of~$f$
and $\Delta_\sfG$ is easily proven to be a continuous group homomorphism from~$\sfG$ to the
multiplicative group~$\real^+$; it is trivial for abelian groups and for compact
groups.

When $\alg=C_{\rm c}(X)$
is the algebra of a commutative dynamical system, the convolution
algebra $C_{\rm c}(\sfG\times X)$ consists of functions $f\colon\sfG\times X\to\complex$ and the
convolution product reads as
\[
(f\star f')(\gamma,x) = \int_\sfG f(\gamma',x)
f'\big(\gamma'^{-1} \gamma,\gamma'^{-1}\cdot x\big)\, \dd \mu_\sfG(\gamma') ,
\]
while the $*$-algebra structure is given by
\[
f^\dag(\gamma,x) = \Delta_\sfG(\gamma)^{-1} \overline{f\big(\gamma^{-1},\gamma^{-1}\cdot x\big)} .
\]
